We quantify the reported Reock proxy's boundary insensitivity by comparing
values computed on smoothed vs.\ raw TIGER/Line district geometries.

\subsection{Boundary Representations}

\paragraph{Raw TIGER/Line.}
The 2020 Census TIGER/Line tract polygons represent coastlines, rivers, and
state boundaries with the full digitisation detail of the source data.
NC coastal tracts may have $> 1{,}000$ vertices along a single coastal segment.

\paragraph{Smoothed boundaries.}
We apply Douglas-Peucker simplification~\citep{douglas1973algorithms}
with tolerance $\epsilon = 100$ metres to reduce vertex counts by approximately
60--80\% for coastal tracts, while preserving the overall shape of inland boundaries.
This represents a simplified-boundary comparison, not a court-certified
``true'' compactness value.

\subsection{Reock vs.\ PP Sensitivity}

Table~\ref{tab:boundary_sensitivity} compares the change in mean district reported Reock
and PP when switching from smoothed to raw boundaries for NC coastal districts
(seed $= 42^\dagger$, ratio-optimal).

\begin{table}[ht]
\centering
\caption{Mean metric change (smoothed $\to$ raw boundaries) for NC coastal
  districts under ratio-optimal, seed $= 42^\dagger$.}
\label{tab:boundary_sensitivity}
\begin{tabular}{lccc}
\toprule
Metric & Smoothed & Raw (TIGER/Line) & Change \\
\midrule
Reported Reock & 0.374$^\dagger$ & 0.355$^\dagger$ & $-0.019$ \\
PP    & 0.261$^\dagger$ & 0.213$^\dagger$ & $-0.048$ \\
\bottomrule
\end{tabular}
\begin{tablenotes}
\footnotesize
\item Note: ``Coastal districts'' = NC districts with $\geq 10$\% of their
  boundary on Atlantic coast or sound (approximately 5 of 14 NC districts).
\item $^\dagger$Single-run result, seed $= 42$.
\end{tablenotes}
\end{table}

The reported Reock decrease ($-0.019$) is 2.5$\times$ smaller than the PP
decrease ($-0.048$).  This supports the operational claim that the current
\textsc{bisect} Reock proxy is less sensitive to boundary representation than PP
for these coastal and riverine districts.  It does not by itself prove an
exact-MBC Reock sensitivity bound.

\subsection{Mechanism}

For exact-MBC Reock, the circle is determined by a small support set of boundary
points, often two or three extremes.  For the current centroid proxy, the radius
is determined by whichever boundary coordinate is farthest from the polygon
centroid.  In both cases, small coastline wiggles usually matter less than they
do for a perimeter sum, but the implemented proxy should be described by its
centroid-radius rule rather than by Welzl support-point mechanics.

In contrast, PP depends on the total perimeter $P = \sum_i |e_i|$ summed over all
boundary edges. Each additional vertex introduced by TIGER/Line digitisation of
the coastline adds perimeter length, directly increasing $P$ and decreasing PP.
